\section{Discussion}
\label{sec:discussion}

The main advantage of the proposed two-scale approach is that it combines detailed stochastic simulation within communities with an aggregate representation of the connections between them. This is suitable for problems in which contact-network data are available at different levels of detail. For example, the internal structure of communities can be reconstructed or generated from known characteristics, whereas inter-community connections may be represented only by aggregate movement flows obtained from transport or migration data. Unlike full-network simulation, \TS{} does not require the nodes involved in inter-community contacts to be specified explicitly. At the same time, unlike fully aggregated models, it preserves individual infection dynamics and the effect of internal topology.

The experiments showed that this simplification can preserve high accuracy in the temporal characteristics of infection spread. In communities with fast internal spread, the arrival-time distributions from \Full{} and \TS{} were close, and this agreement persisted even after nine consecutive inter-community crossings. Thus, the absence of information about specific bridge nodes does not necessarily cause a substantial error in the overall propagation dynamics. This makes aggregate interfaces a promising way to model networks of communities when the main quantities of interest are the times at which communities are reached and their distributions.

The accuracy of this approximation can be explained by the ratio of time scales. If infection spreads rapidly among the nodes of a community relative to the waiting time for inter-community transmission, the specific locations of the bridges have a comparatively small effect on the export time. With slower and more heterogeneous internal spread, however, the states of individual bridge nodes may remain important for a substantial part of the waiting time for a crossing. This is why structural characteristics, including ASPL, do not fully explain the observed discrepancies. The stronger relationship between accuracy and the measure $\rho=\beta_{\mathrm{out}}\sigma_T$ supports an explanation based on temporal heterogeneity in internal infection. This suggests that the suitability of aggregation can be assessed from both network topology and process dynamics.

A possible direction for further research is to increase the model's level of detail by incorporating local community characteristics, such as ASPL, eigenvalues, or other descriptors, and to evaluate whether these additions reduce the discrepancy between the methods.

For simulations of spread through many communities, the way in which discrepancies accumulate is also important. The experiments showed that the absolute difference between arrival times can increase along the chain, whereas the relative shift often stabilises or changes slowly. Thus, the error of the aggregate simulation does not necessarily grow in proportion to the number of inter-community crossings. At the same time, the stronger bias in the total duration of later crossings than in the first crossing indicates possible additional effects after infection is imported. One explanation is the occurrence of surplus imports, which can create additional foci of infection and accelerate subsequent spread. The separate contribution of this mechanism requires further investigation.

The theoretical analysis clearly identifies the approximation implemented by the two-scale approach. The ideal process is equivalent to exact simulation on a surrogate network with uniformly distributed inter-community connections. This provides a mathematical basis for the aggregate representation and explains why its results can differ from those of full-network simulation. In particular, similarity between time distributions does not guarantee reproduction of specific infection routes. The appropriate level of detail should therefore be determined by the aim of the study: an aggregate representation may be sufficient for estimating temporal characteristics, whereas analysis of physical transmission paths requires information about the corresponding contacts.

Further possibilities arise from the computational structure of the method. Independent local simulators can run in parallel between synchronisation points, and individual communities can have different internal structures and simulation parameters. In addition, replacing explicit inter-community edges with aggregate weights can reduce the amount of contact data. However, the actual gains in speed and memory depend on the network structure and the cost of synchronisation.

\section{Conclusions}
\label{sec:conclusions}

This study investigated a two-scale method for stochastic simulation of infection spread in networks of communities. The method combines detailed Gillespie simulation of internal contact networks with aggregate inter-community intensities. It therefore allows local contact networks to be used together with aggregate data on connections between communities, without requiring every inter-community contact to be represented explicitly. Its modular structure also allows local simulators to be run in parallel.

We proved that the ideal two-scale process is equivalent to exact simulation on a surrogate network with uniform inter-community connections. We analysed the mechanisms that produce discrepancies between aggregate and full-network simulation, including the consequences of losing information about specific bridge nodes. We also proved that, when bridge nodes are placed uniformly at random, aggregation yields a first infection export that is stochastically no later in the ideal model.

The experiments confirmed that \TS{} can reproduce infection-spread time distributions close to those of \Full{} when infection spreads rapidly within communities. This agreement persisted even after nine consecutive inter-community crossings. The absolute discrepancy between arrival times can accumulate along a chain, whereas the relative discrepancy often stabilises or changes slowly.

We examined how internal structure, community size, and the number of inter-community edges affect the accuracy of aggregate simulation. The time-scale ratio $\rho=\beta_{\mathrm{out}}\sigma_T$ had a stronger observed relationship with the discrepancies than ASPL and is a promising characteristic for identifying conditions under which the method is applicable.

In our tests, the smallest discrepancies occurred when infection spread rapidly within communities. In all 18 configurations with $\rho<0.25$, the normalised Wasserstein distance did not exceed 4.5\%. At the same time, close time distributions did not guarantee preservation of admissible transmission routes.

\clearpage
